\pdfoutput=1

\documentclass[11pt]{article}

\usepackage{EACL2023}
\usepackage{times}
\usepackage{latexsym}
\usepackage[T1]{fontenc}
\usepackage[utf8]{inputenc}
% natbib and setcitestyle lines REMOVED
\usepackage{microtype}
\usepackage{inconsolata}
\usepackage{graphicx}
\usepackage{float}
\usepackage{placeins}
\usepackage{xcolor}
\usepackage{booktabs}
\usepackage{enumitem}


\PassOptionsToPackage{square,numbers}{natbib}


% add this too
% ---------------------------------------------------

\title{Multi-Horizon Stock Price Trajectory Forecasting:
A Hybrid Framework Integrating Contextual and
Sectoral Sentiment Analysis}

\author{
    \begin{minipage}{0.48\textwidth}
    \textbf{Team Name: The Tokenizers} \\[0.5em]
    \small
    \begin{tabular}{c l c}
        \hline
        \textbf{S.No} & \textbf{Name} & \textbf{Roll No.} \\
        \hline
        1 & Naveen Mishra & 2025201084 \\
        2 & C. Phani Rama Vaibhav & 2025201067 \\
        3 & Akula Amanaganti & 2025202034 \\
        4 & Sriram Paruchuri & 2025201074 \\
        \hline
    \end{tabular}
    \end{minipage}
}

\setlength\titlebox{5cm}

\begin{document}
\maketitle

\begin{abstract}
Stock prices are influenced by a complex ecosystem of financial news, market sentiment, and geopolitical events. However, traditional prediction models often exhibit a form of “tunnel vision” by relying primarily on historical price movements or shallow sentiment signals. Such approaches do not capture latent dependencies in the financial ecosystem, for example, how disruptions affecting a Tier-1 supplier may indirectly impact the stock price of a manufacturer after a delay.

To address this limitation, we propose a hybrid framework that integrates contextual sentiment analysis with sectoral information to capture deeper relationships between news events and stock price trajectories. By modeling these hidden signals, the proposed approach aims to improve multi-horizon stock price forecasting and provide a more comprehensive understanding of market dynamics.
We propose a context-sensitive hybrid framework that combines historical OHLC data and news sentiment to forecast percentage changes in stock price over a future window of Y days using an observation window of X days.


\end{abstract}

\textbf{Keywords:} Stock Price Forecasting, Financial Sentiment Analysis, FinBERT, LSTM, Temporal Fusion Transformer, Time Series Prediction, News-Based Market Analysis, RMSE Optimization

\section{Introduction}

Financial markets are highly dynamic systems influenced by a wide range of factors including macroeconomic indicators, geopolitical developments, industry-specific events, and investor sentiment. In recent years, machine learning and deep learning techniques have been widely applied to stock price forecasting due to their ability to model complex nonlinear relationships in financial data. Traditional time-series approaches primarily rely on historical price indicators such as Open, High, Low, and Close (OHLC) values, while more recent approaches incorporate textual sentiment extracted from financial news sources.

Despite these advancements, many existing models exhibit a limited understanding of the broader financial ecosystem. Most prediction systems focus on direct relationships between news sentiment and individual stock prices without considering indirect dependencies across sectors and supply chains. For example, disruptions affecting a major supplier may influence the performance of downstream manufacturers after a delay, a relationship that conventional models often fail to capture. This limitation results in incomplete modeling of market dynamics and reduces forecasting accuracy.

To address these challenges, this work proposes a context-sensitive hybrid framework that integrates historical stock price data with contextual sentiment extracted from financial news. The framework aims to capture hidden relationships between market events and stock price movements by combining time-series forecasting with domain-specific natural language processing techniques.

The proposed methodology collects historical stock data from 2015 onward along with large-scale textual information extracted from financial news, business reports, and general news articles. Sentiment representations are generated using a financial-domain transformer model and integrated with time-series data to predict percentage changes in stock prices over a future forecasting window.

Several experimental configurations are implemented to evaluate the effectiveness of different observation windows and model architectures. The primary focus of this study is on Indian NIFTY 50 stocks, where stock price data is combined with sentiment signals generated from financial news articles, business reports, and related textual sources. Experiments are conducted using six-month and three-year stock datasets with multiple observation windows of 7, 10, 15, and 30 days, employing FinBERT combined with Long Short-Term Memory (LSTM) networks. 

Additional experiments are performed on a six-month dataset using a FinBERT and Temporal Fusion Transformer (TFT) architecture. For comparative and learning purposes, large-scale experiments are also conducted using NASDAQ stock data from 1999–2023 combined with sentiment signals. The performance of these configurations is evaluated using Mean Absolute Percentage Error (MAPE) and Root Mean Square Error (RMSE) to identify the most effective model architecture and observation window for forecasting.

The long-term objective of this research is to develop a predictive system capable of dynamically adapting to incoming financial data while providing accurate multi-horizon stock price trajectory predictions under evolving market conditions.

\section{Related Work}
Financial sentiment analysis using FinBERT with application in predicting stock movement (T. Jiang et al., 2025) \cite{T.Jiang}{} proposed a hybrid deep learning framework that integrates FinBERT-based financial sentiment analysis with Long Short-Term Memory (LSTM) networks to improve stock market prediction. In their approach, financial news articles are processed using FinBERT to extract sentiment probabilities (positive, negative, and neutral), which are aggregated into a daily sentiment score. These sentiment scores are then combined with historical stock market features such as open, high, low, close prices and trading volume, and used as input to an LSTM model for time-series forecasting of stock prices. The study compares this model with baseline approaches including ARIMA, standalone LSTM, BERT + LSTM, and FinBERT + DNN, and demonstrates that the FinBERT + LSTM model achieves the lowest prediction error, indicating that incorporating financial sentiment information significantly improves stock movement prediction accuracy.


A Novel Ensemble Deep Learning Model for Stock Prediction Based on Stock Prices and News (Y. Li et al., 2021) \cite{Y.Li} proposes a deep learning framework that combines financial news sentiment with historical stock price data using a blending ensemble architecture. The study notes that stock prices are influenced by factors such as company news, company performance, industry trends, investor sentiment and economic conditions, and emphasizes the importance of selecting an appropriate historical window size for prediction. Financial news from sources such as CNBC, Reuters, Wall Street Journal and Fortune are collected, and sentiment scores are extracted from news titles using the VADER sentiment analysis tool. The dataset is divided into training, validation and test sets, and a rolling window of 10 days is used to capture temporal dependencies between news sentiment and stock prices. The proposed model uses a two-level blending ensemble architecture where the first level consists of two recurrent neural networks, LSTM and GRU, which learn temporal patterns from combined sentiment scores and stock price data. Each model contains four layers with 50 neurons and a dropout rate of 0.2 to reduce overfitting. Predictions from these models on the validation dataset are used as inputs to a second-level meta learner implemented as a fully connected neural network with three layers and ReLU activation. The model is evaluated using metrics such as Mean Squared Error (MSE), Mean Prediction Accuracy (MPA), precision, recall, F1-score and Movement Direction Accuracy (MDA). Experimental results on S\&P 500 data show that the blending ensemble model outperforms individual models such as LSTM, GRU and DP-LSTM, reducing MSE by up to 57.55\% and achieving about 66.67\% movement direction accuracy.

Stock Price Prediction Using News Sentiment Analysis (S. Mohan et al., 2019) \cite{S.Mohan} proposed a combination of models to predict stock price movements by incorporating financial news sentiment analysis. Their work predicts stock price changes by analyzing the polarity of financial news, indicating potential increases or decreases in stock value. The study evaluates several models including ARIMA, Facebook Prophet, and RNN-LSTM using a sliding window approach on historical stock price data. 

The authors collected five years of S\&P 500 stock price data along with more than 265,000 financial news articles, enabling analysis of the relationship between news sentiment and stock price movements. Textual data was preprocessed using techniques such as HTML tag removal, tokenization, and feature extraction, and sentiment polarity was extracted using NLTK-based sentiment analysis.

Different experimental setups compared models using only historical price data and models incorporating sentiment polarity from financial news. Among the approaches explored, the RNN-pp model (RNN-LSTM with prices and text polarity as input) achieved the best performance. The models were evaluated using Mean Absolute Percentage Error (MAPE), where deep learning approaches showed lower error compared to traditional models such as ARIMA and Facebook Prophet.



\section{Dataset Details}

This study utilizes financial news and stock market data related to companies listed in the NIFTY-50 index. The dataset integrates historical stock price information with sentiment extracted from financial news articles to analyze the relationship between news sentiment and stock price movements.

\subsection{Indian NIFTY-50 Dataset}

The dataset covers stock market data from January 2023 to march 2026 for multiple companies in the NIFTY-50 index. \textbf{Training Strategy:} The overall dataset spans the period from 2015 to 2026 and contains several thousand daily trading records after preprocessing. The primary objective of this work is to develop a robust stock prediction model capable of effectively utilizing this large historical dataset. As an initial step toward this goal, experiments are currently conducted using shorter observation horizons of six months and three years. These smaller subsets are used to systematically evaluate suitable window sizes, hyperparameters, and model configurations before scaling the approach to the complete dataset.Financial news articles are collected and aligned with the corresponding stock symbols and trading dates. Each record combines stock price indicators with processed news sentiment scores.

The main columns used in the dataset are:

\begin{itemize}
\item \textbf{Date} – Trading date corresponding to the stock price record.
\item \textbf{Ticker} – Stock symbol representing the company listed in the NIFTY-50 index.
\item \textbf{Open} – Opening price of the stock at the beginning of the trading day.
\item \textbf{High} – Highest price reached by the stock during the trading day.
\item \textbf{Low} – Lowest price reached during the trading day.
\item \textbf{Close} – Final trading price of the stock at market close.
\item \textbf{Adj Close} – Adjusted closing price accounting for dividends and corporate actions.
\item \textbf{Volume} – Total number of shares traded during the day.
\item \textbf{date} – Timestamp corresponding to the news article publication.
\item \textbf{cleaned\_news} – Preprocessed financial news text after removing noise, punctuation and stop words.
\item \textbf{cleaned\_news\_len} – Length of the cleaned news text used as a textual feature.
\item \textbf{neg\_score} – Negative sentiment score extracted from the news text.
\item \textbf{neu\_score} – Neutral sentiment score extracted from the news text.
\item \textbf{pos\_score} – Positive sentiment score extracted from the news text.
\end{itemize}

The dataset is prepared in two configurations to evaluate the impact of different observation windows :

\textbf{Six-Month Dataset}
\begin{itemize}
\item Total news articles: 4046
\item Merged dataset size: (5586, 95)
\item Training data shape: (4459, 95)
\item Validation data shape: (539, 95)
\item Test data shape: (588, 95)
\end{itemize}

\textbf{Three-Year Dataset}
\begin{itemize}
\item Total news articles: 30,663
\item Merged dataset size: (38,181, 14)
\item Training data shape: (30,005, 97)
\item Validation data shape: (3,779, 97)
\item Test data shape: (3,847, 97)
\end{itemize}

\subsection{Additional NASDAQ Financial News Dataset}
In addition to the Indian dataset, a large-scale financial news dataset related to NASDAQ-listed companies was used during earlier stages of model training. This dataset consists of approximately 15.7 million financial news articles and around 29.7 million stock price records covering about 4,775 companies from the S\&P 500 index over the period from 1999 to 2023. The dataset includes information such as article titles, publication dates, publishers, authors, article content, and generated summaries. The large scale and long historical coverage of this dataset allow the model to learn general financial language patterns and sentiment representations. Pretraining on this dataset helps the model better generalize and adapt when applied to the NIFTY-50 dataset containing Indian financial news and stock price data.

\section{Data Preprocessing}

Before training the model, several preprocessing steps were applied to ensure compatibility between textual sentiment features and numerical financial indicators. Missing entries were removed and the dataset was sorted chronologically to maintain temporal consistency.

All numerical features were normalized using Min-Max scaling in order to improve neural network training stability. Min-Max normalization rescales each feature to the interval $[0,1]$ according to the transformation:

\[
x_{scaled} = \frac{x - x_{min}}{x_{max} - x_{min}}
\]

This normalization was applied to all financial indicators including open, high, low, close and volume as well as the sentiment scores extracted from financial news.

Because the dataset contains multiple stock tickers, each ticker symbol was converted into a numerical identifier using enumeration. Each unique ticker was assigned an integer value which was subsequently scaled using Min-Max normalization. The scaled ticker identifier was included as a static feature in the model input so that the neural network can distinguish between different stocks during training.

After normalization the dataset was transformed into sequential training samples using a sliding window approach. For each prediction step the model observes the previous $X$ days of data and predicts the next day closing price. This transformation converts the dataset into a three dimensional tensor representing samples, time steps and features. In the best performing configuration the input tensor has the shape $(batch,11,9)$ where 11 represents the number of time steps and 9 represents the total number of input features.

\section{Methodology}

\subsection{Preprocessing and Sentiment Score Calculation}

The collected financial news data undergoes several preprocessing steps before sentiment extraction. First, the raw news articles are cleaned by removing HTML tags, URLs, and other formatting artifacts. Articles containing multiple line breaks are converted into a single continuous line, and all text is transformed to lowercase to ensure consistency during processing.

After cleaning, the news articles are aggregated based on their publication date so that all articles published on the same day are combined into a single text representation. This aggregated daily news text is then passed to the FinBERT model to compute sentiment scores. FinBERT produces three sentiment probabilities for each input: positive, negative, and neutral.

The resulting daily sentiment scores are merged with the corresponding NIFTY-50 stock price data using a left join based on the date and ticker symbol. This ensures that each trading record is associated with the sentiment information available for that day. Since the same global sentiment is applied to all companies for a given date, the dataset expands approximately fifty times to match each ticker symbol. The merged dataset is then stored for further processing in the modeling stage.


\subsection{Model Creation and Hyperparameter Tuning}

In the model training phase, the merged dataset is first loaded and prepared for training. Numerical stock price features such as open, high, low, close, adjusted close, and volume are normalized using Min-Max normalization to scale the values into the range of 0 to 1. The ticker symbols are converted into numerical representations and similarly scaled.

The dataset is then divided into training, validation, and testing sets based on the selected observation window sizes. Multiple window sizes are evaluated, specifically 7, 10, 15, and 30 days, in order to analyze how different temporal contexts influence the prediction performance.

For each window size, hyperparameter tuning is performed to identify the most suitable model configuration. The parameters tuned include the learning rate, the number of LSTM units in each of the two LSTM layers, dropout rate to prevent overfitting, and the number of units in the dense layer. After training models with different parameter combinations, the performance of each configuration is evaluated using the test dataset. The best performing model is selected based on its predictive accuracy and error metrics.


\subsection{Evaluation Metrics}

Model performance was evaluated using several regression metrics including Mean Squared Error, Mean Absolute Error, Root Mean Squared Error and Mean Absolute Percentage Error.

\[
MSE = \frac{1}{n}\sum (y - \hat{y})^2
\]

\[
MAE = \frac{1}{n}\sum |y - \hat{y}|
\]

\[
RMSE = \sqrt{MSE}
\]

\[
MAPE = \frac{100}{n}\sum \left|\frac{y - \hat{y}}{y}\right|
\]

\begin{figure}[htbp]
\centering
\includegraphics[width=0.85\linewidth, keepaspectratio]{architecture.jpg}
\caption{Proposed Architecture for Stock Price Prediction Framework}
\label{fig:architecture}
\end{figure}

\begin{figure}[htbp]
\centering
\includegraphics[width=0.75\linewidth, keepaspectratio]{tft.jpg}
\caption{Proposed Architecture for FinBERT + Temporal Fusion Transformer (TFT)}
\label{fig:tft_architecture}
\end{figure}

\section{Results}
In this section we present the results obtained for hyperparameter tuning, the evaluation metrics observed observed on test set for best model on each window size, and the summary of the best model obtained so far.

\subsection{Best Model Architecture}

The best performing model configuration was obtained using an observation window of 10 days. The network was implemented using the Keras Functional API and contains two LSTM layers with dropout regularization followed by dense layers.

\begin{table}[h]
\centering
\small
\begin{tabular}{l r r}
\hline
\textbf{Layer} & \textbf{Output Shape} & \textbf{Parameters} \\
\hline
Input Layer & (None, 11, 9) & 0 \\
LSTM Layer 1 & (None, 11, 64) & 18,944 \\
Dropout & (None, 11, 64) & 0 \\
LSTM Layer 2 & (None, 32) & 12,416 \\
Dropout & (None, 32) & 0 \\
Dense Layer & (None, 64) & 2,112 \\
Output Layer & (None, 1) & 65 \\
\hline
\textbf{Total} & & \textbf{33,537} \\
\hline
\end{tabular}
\caption{Architecture of the best performing LSTM model}
\end{table}

The total number of trainable parameters in the model is 33,537.


\subsection{Evaluation Results}
The results show that the 10 day observation window provides the best predictive performance across all evaluation metrics. Shorter windows fail to capture sufficient historical context while longer windows introduce additional noise that reduces prediction accuracy.

\begin{table}[h]
\centering
\begin{tabular}{ccccc}
\hline
Window & MSE & MAE & RMSE & MAPE \\
(days) & & & & (\%) \\
\hline
7 & $2.32\times10^5$ & 126.64 & 482.17 & 4.21 \\
\textbf{10} & $\mathbf{7.35\times10^3}$ & \textbf{43.70} & \textbf{85.71} & \textbf{3.28} \\
15 & $9.64\times10^6$ & 1825.83 & 3104.23 & 177.50 \\
30 & $2.80\times10^5$ & 153.45 & 528.98 & 6.24 \\
\hline
\end{tabular}
\caption{Performance comparison across different window sizes}
\end{table}

The results show that the 10 day observation window provides the best predictive performance across all evaluation metrics. Shorter windows fail to capture sufficient historical context while longer windows introduce additional noise that reduces prediction accuracy.

\subsection{Visual Results}

\begin{figure}[H]
\centering
\includegraphics[width=0.65\linewidth, height=0.5\textheight, keepaspectratio]{lowest_mape.png}
\caption{Top 5 Stocks with Lowest MAPE --- Actual vs Predicted Close Prices.}
\label{fig:lowest_mape}
\end{figure}

\begin{figure}[H]
\centering
\includegraphics[width=0.85\linewidth, height=\textheight]{windowise_comparision.jpeg}
\caption{Evaluation Metrics by Window Size across 7, 10, 15, and 30-day observation windows.}
\label{fig:windowwise}
\end{figure}



\begin{figure}[H]
\centering
\includegraphics[width=\linewidth, keepaspectratio]{lstm_actual_vs_predicted.png}
\caption{LSTM Model: Actual vs Predicted Close Prices on Test Set.}
\label{fig:lstm_overall}
\end{figure}


\section{Future Work}

We identify four extensions to the current LSTM + FinBERT pipeline, targeting completion by \textbf{April 8, 2026}:

\begin{enumerate}[noitemsep, topsep=3pt, leftmargin=*]
  \item \textbf{Four-Tier Sentiment Segregation:} Currently every stock receives an identical aggregate sentiment score. We will use NER to route each article into one of four tiers: (1)~Direct (company-specific), (2)~Sectoral (industry trends), (3)~Supply-Chain (supplier events as early-warning signals), and (4)~Geopolitical (macro events). FinBERT scores each tier independently, producing a \textbf{4-dimensional sentiment vector} per company per day.

  \item \textbf{Price-Difference ($\Delta P_t$) Prediction:} Rather than predicting the non-stationary raw closing price, we reformulate the target as $\Delta P_t = P_t - P_{t-1}$. This improves stationarity, enables cross-company generalisation, and directly quantifies directional movement. Evaluation uses RMSE, MAPE, and Directional Accuracy.

  \item \textbf{Temporal Fusion Transformer (TFT):} We will implement a TFT model which offers Variable Selection Networks (learns feature importance per timestep), Gated Residual Networks (stabilises training), multi-head self-attention (captures long-range dependencies), and quantile outputs (P10/P50/P90 probabilistic forecasts).

  \item \textbf{Ensemble Learning:} We will build a stacking ensemble combining LSTM (sequential memory), GRU (lightweight, gradient-stable), and TFT (attention-based, probabilistic) with a learned weighted combiner tuned on validation data.
\end{enumerate}

\section{Timeline}

\begin{table}[h]
\centering
\small
\begin{tabular}{c l c l}
\hline
\textbf{Phase} & \textbf{Task} & \textbf{Duration} & \textbf{Milestone} \\
\hline
1 & Sentiment Segregation & Mar 8--21 & 4-tier dataset ready \\
2 & $\Delta P_t$ Prediction & Mar 15--25 & LSTM-$\Delta$ baseline \\
3 & TFT Implementation & Mar 22--Apr 2 & TFT trained \& evaluated \\
4 & Ensemble Learning & Mar 29--Apr 6 & Ensemble results final \\
5 & Report \& Submission & Apr 5--8 & \textbf{Deadline: Apr 8} \\
\hline
\end{tabular}
\caption{Project timeline for future work (March 8 -- April 8, 2026)}
\end{table}
\section{References}
\subsection{GitHub Link and Dataset Link}
\begin{itemize}
  \item Git-Link: https://github.com/Rama-Vaibhav/Stock-Price-Prediction-using-INLP
  \item Dataset Link:https://tinyurl.com/ms82yft2
\end{itemize}

\begin{thebibliography}{9}

\bibitem{T.Jiang}
T. Jiang and Q. Zeng, 
``Financial Sentiment Analysis Using FinBERT with Application in Predicting Stock Movement,'' 
\textit{arXiv preprint arXiv:2306.02136}, 2025. 
Available: https://arxiv.org/abs/2306.02136.

\bibitem{Y.Li}
Y. Li and Y. Pan, 
``A Novel Ensemble Deep Learning Model for Stock Prediction Based on Stock Prices and News,'' 
\textit{arXiv preprint arXiv:2007.12620}, 2020. 
Available: https://arxiv.org/abs/2007.12620.

\bibitem{S.Mohan}
S. Mohan, S. Mullapudi, S. Sammeta, P. Vijayvergia, and D. C. Anastasiu, 
``Stock Price Prediction Using News Sentiment Analysis,'' 
in \textit{2019 IEEE Fifth International Conference on Big Data Computing Service and Applications (BigDataService)}, 
Newark, CA, USA, 2019, pp. 205--208. 
doi: 10.1109/BigDataService.2019.00035.

\end{thebibliography}



\end{document}